\documentclass[../../main.tex]{subfiles}
\begin{document}
\begin{oframed}
  \subfile{problem}
\end{oframed}

{\bf [解答]}

\paragraph{(1)}

排反事象を考える．
$X$が3でわり切れないのは，$n$回とも1, 2, 4, 5のいずれかが出るときである．
従って求める確率は
\begin{align}
  1 - \left(\dfrac{4}{6}\right)^n = 1 - \left(\dfrac{2}{3}\right)^n \label{eq:1}
\end{align}
である．

\paragraph{(2)}

\begin{figure}[htbp]
  \centering
  \begin{tikzpicture}[x=1cm, y=1cm, >=stealth, every node/.style={font=\small}]
    % 全体集合 U の枠
    \draw[line width=1.0pt, gray!70!black, rounded corners=6pt] (-3.8, -2.1) rectangle (3.8, 2.1);
    \node[font=\bfseries, gray!80!black, below right=3pt] at (-3.8, 2.1) {$U$};

    % 円のパラメータ
    \def\r{1.6}
    \coordinate (cA) at (-0.9, 0);
    \coordinate (cB) at (0.9, 0);

    % 集合 A の塗り
    \fill[blue!15, opacity=0.6] (cA) circle (\r);
    % 集合 B の塗り
    \fill[teal!15, opacity=0.6] (cB) circle (\r);

    % 積集合 A ∩ B の塗り
    \begin{scope}
      \clip (cA) circle (\r);
      \fill[blue!40!teal!35, opacity=0.85] (cB) circle (\r);
    \end{scope}

    % 円の境界線
    \draw[line width=1.2pt, blue!85!black] (cA) circle (\r);
    \draw[line width=1.2pt, teal!85!black] (cB) circle (\r);

    % ラベル
    \node[blue!90!black, font=\bfseries\normalsize] at (-1.75, 0.25) {$A$};
    \node[blue!80!black, font=\scriptsize] at (-1.75, -0.25) {$3$の倍数};

    \node[teal!90!black, font=\bfseries\normalsize] at (1.75, 0.25) {$B$};
    \node[teal!80!black, font=\scriptsize] at (1.75, -0.25) {$2$の倍数};

    \node[purple!75!black, font=\bfseries\normalsize] at (0, 0.25) {$A \cap B$};
    \node[purple!75!black, font=\scriptsize] at (0, -0.25) {$6$の倍数};
  \end{tikzpicture}
  \caption{事象 $A, B$ のベン図と積事象 $A \cap B$}
  \label{fig:venn_ab}
\end{figure}

事象を
\begin{itemize}
  \item A: $X$が3でわり切れる \\
  \item B: $X$が2でわり切れる
\end{itemize}
と定めると，「$X$が6で割り切れる」事象は$A \cap B$に等しいから，求める確率は
\begin{align}
  P(A \cap B) = P(A) + P(B) - P(A \cup B) \label{eq:2}
\end{align}
である．$P(A)$は\cref{eq:1}であり，$B, A \cup B$も(1)と同様に排反事象を考えて
\begin{align}
  P(B) &= 1 - \left(\dfrac{3}{6}\right)^n = 1 - \left(\dfrac{1}{2}\right)^n \\
  P(A \cup B) &= 1 - \left(\dfrac{2}{6}\right)^n = 1 - \left(\dfrac{1}{3}\right)^n
\end{align}
だから，\cref{eq:2}に代入して求める確率は
\begin{align}
  P(A \cap B)
  &= 1 - \left(\dfrac{2}{3}\right)^n + 1 - \left(\dfrac{1}{2}\right)^n - \left\{1 - \left(\dfrac{1}{3}\right)^n\right\} \\
  &= 1 + \left(\dfrac{1}{3}\right)^n - \left(\dfrac{1}{2}\right)^n - \left(\dfrac{2}{3}\right)^n
\end{align}
である．

\end{document}
